\documentclass[a4paper,11pt]{article}
\usepackage[T1]{fontenc}
\usepackage[utf8]{inputenc}
\usepackage[left=2cm,right=2cm,top=2.5cm,bottom=2.5cm]{geometry}
\setlength{\headheight}{14pt}
\usepackage{xcolor}
\usepackage{mdframed}
\usepackage{parskip}
\usepackage{fancyhdr}
\usepackage{hyperref}

\definecolor{usercolor}{RGB}{30, 100, 200}
\definecolor{agentcolor}{RGB}{30, 150, 80}
\definecolor{doccolor}{RGB}{200, 90, 10}
\definecolor{userbg}{RGB}{235, 244, 255}
\definecolor{agentbg}{RGB}{235, 250, 238}
\definecolor{docbg}{RGB}{255, 243, 220}
\definecolor{answerbg}{RGB}{250, 250, 235}

\newmdenv[backgroundcolor=userbg,linecolor=usercolor,linewidth=1.5pt,
  leftmargin=0pt,rightmargin=80pt,innerleftmargin=10pt,innerrightmargin=10pt,
  innertopmargin=8pt,innerbottommargin=8pt,roundcorner=5pt,
  skipabove=6pt,skipbelow=3pt]{userbubble}

\newmdenv[backgroundcolor=agentbg,linecolor=agentcolor,linewidth=1.5pt,
  leftmargin=80pt,rightmargin=0pt,innerleftmargin=10pt,innerrightmargin=10pt,
  innertopmargin=8pt,innerbottommargin=8pt,roundcorner=5pt,
  skipabove=3pt,skipbelow=3pt]{agentbubble}

\newmdenv[backgroundcolor=docbg,linecolor=doccolor,linewidth=1pt,
  leftmargin=80pt,rightmargin=0pt,innerleftmargin=10pt,innerrightmargin=10pt,
  innertopmargin=6pt,innerbottommargin=6pt,roundcorner=4pt,
  skipabove=2pt,skipbelow=3pt]{docenv}

\newmdenv[backgroundcolor=answerbg,linecolor=gray,linewidth=0.5pt,
  leftmargin=80pt,rightmargin=0pt,innerleftmargin=10pt,innerrightmargin=10pt,
  innertopmargin=5pt,innerbottommargin=5pt,roundcorner=4pt,
  skipabove=2pt,skipbelow=6pt]{refanswerenv}

\pagestyle{fancy}
\fancyhf{}
\renewcommand{\headrulewidth}{0.3pt}
\cfoot{\thepage}

\lhead{\textcolor{gray}{\small RAG Conversation Dataset}}
\rhead{\textcolor{gray}{\small Conv.\,\#2\ --\ govt0611}}
\begin{document}

\begin{center}
{\LARGE\bfseries RAG Conversation #3}
\par\medskip
{\small\textcolor{gray}{Domain:\ \texttt{mt-rag-govt-elser-512-100-20240611}\quad|\quad Status:\ \textbf{accepted}}}
\end{center}
\vspace{0.4em}
\noindent\textcolor{gray}{\rule{\linewidth}{0.6pt}}
\vspace{0.4em}
\vspace{0.5em}
\noindent{\small\textcolor{usercolor}{\textbf{$\triangleright$\enspace User --- Turn~1}}}\hfill{\tiny\textcolor{gray}{Summarization $\cdot$ ANSWERABLE}}
\begin{userbubble}
What is ground water contamination?
\end{userbubble}
\noindent\hfill{\small\textcolor{agentcolor}{\textbf{Agent\enspace$\triangleleft$}}}
\begin{agentbubble}
Groundwater contamination refers to the presence of harmful substances, such as household cleaners, flea powders, garden pesticides, paints, varnishes, motor oil, and antifreeze. It occurs when these products are disposed of improperly and seep into the ground.  By disposing of these harmful products improperly, we are contributing to groundwater contamination.
\end{agentbubble}
\noindent\hfill{\footnotesize\textcolor{doccolor}{\textbf{$\hookrightarrow$\enspace Retrieved documents~(2)}}}
\begin{docenv}
{\footnotesize\textcolor{doccolor}{\textbf{Doc~1}}}\hfill{\tiny\texttt{score:~13.313}}\ {\tiny\textcolor{gray}{\{'relevant': \{'9dcd142c-415e-4161-a82f-c8b14bfb36ac': \{'value': 'yes', 'timestamp': 1725438894\}\}\}}}\\[3pt]
{\small\textit{FAQs • Monterey Park, CA • CivicEngage}}\\[2pt]
{\tiny\url{https://www.montereypark.ca.gov/faq.aspx}}\\[4pt]
{\footnotesize FAQs • Monterey Park, CA • CivicEngage
Rain and snowmelt comes down from the San Gabriel Mountains through a series of dams and are then diverted to "spreading basins". Spreading basins are designated areas of land that are very permeable, such as sand and gravel areas. Water can then percolate into the ground. 100\% of the water served to residents and businesses in Monterey Park comes from groundwater. Thus, keeping our groundwater resource clean is vitally important to all of us. 	

7.
Where do groundwater toxins come from?

Many products we use around our homes everyday are considered toxic, such as household cleaners, flea powders, garden pesticides, paints, varnishes, motor oil, and antifreeze. These products are not only hazardous to our health, but to our environment. By disposing these harmful products improperly, we are contributing to the groundwater contamination in the San Gabriel Valley.	

8.
Can contaminated groundwater be cleaned up?

Yes, but it can be time-consuming and expensive. Communities whose drinking water source has become contaminated have had to spend millions of dollars to remove contaminants from the water before it can be piped to homes and businesses. The cleanup is effective, but it can significantly affect the cost of providing the water to customers. It is far better to prevent contamination in the first place. What are the dos and do nots for groundwater?Do: Purchase only what you need Use products completely before disposing Recycle whatever you can (oil, batteries, etc.) Recycle surplus paintsDo not: Dump products onto the ground (they can contaminate soil and groundwater) Pour products or wastes into the street (they can contaminate oceans, lakes, and groundwater) Put hazardous or toxic products in the trash Pour products down the drainView proper disposal of Household Hazardous Wastes online. 

9.
Is our water fluoridated?

At this time fluoride is not added to the water. There are low levels of fluoride naturally occurring in the water. The range at which they occur is between 0.2 to 0.8 milligrams per liter (mg/L).}
\end{docenv}
\begin{docenv}
{\footnotesize\textcolor{doccolor}{\textbf{Doc~2}}}\hfill{\tiny\texttt{score:~20.957}}\ {\tiny\textcolor{gray}{\{'relevant': \{'9dcd142c-415e-4161-a82f-c8b14bfb36ac': \{'value': 'yes', 'timestamp': 1725439125\}\}\}}}\\[3pt]
{\small\textit{FAQs • Where do groundwater toxins come from?}}\\[2pt]
{\tiny\url{https://www.montereypark.ca.gov/FAQ.aspx?QID=118}}\\[4pt]
{\footnotesize FAQs • Where do groundwater toxins come from?
FAQs • Where do groundwater toxins come from?

 

Skip to Main Content

Create a Website Account - Manage notification subscriptions, save form progress and more.   

Website Sign In

 

 中文 Español

 
 
 

GovernmentServicesResidentsVisitorsBusinessHow Do I...

 

 

 

 

 

 

 

 

 
 
 

 

Search

 
 
 

 
HomeFAQs

Search

All categories
Code Enforcement
Construction Inspections
Drone FAQs
Election Information (English)
Fire Department
Library
Multi-Family Residential/Apartment Inspection Program
Planning
Public Works - Slurry Seal
Water

Categories

All Categories

Code Enforcement

Construction Inspections

Drone FAQs

Election Information (English)

Fire Department

Library

Multi-Family Residential/Apartment Inspection Program

Planning

Public Works - Slurry Seal

Water

 

Where do groundwater toxins come from?	

Many products we use around our homes everyday are considered toxic, such as household cleaners, flea powders, garden pesticides, paints, varnishes, motor oil, and antifreeze. These products are not only hazardous to our health, but to our environment. By disposing these harmful products improperly, we are contributing to the groundwater contamination in the San Gabriel Valley.	

▼
Water	

Show All Answers

1.
How do I pay my water bill?

Water bills are mailed out on a two-month cycle. There are three ways to pay your bill: by mail, in person, or automatic billing.Bills can be paid in person at the Cashiers Counter located in the first floor of the Monterey Park City Hall at 320 W. Newmark Ave. View the location map online. The city has recently established a new automated billing system where your water bill can be automatically paid from your checking account. See more information on automated billing. To mail a payment, send it to Water / Refuse Billing at 320 West Newmark Avenue, Monterey Park, CA 91754. For additional information about your water bill you may call 626-307-1342. 

2.
How do I report problems such as dirty water or a water leak in the street?}
\end{docenv}
\begin{refanswerenv}
{\footnotesize\textit{\textcolor{gray}{\textbf{Reference answer:}\enspace Groundwater contamination refers to the presence of harmful substances, such as household cleaners, flea powders, garden pesticides, paints, varnishes, motor oil, and antifreeze, in the groundwater. This occurs when these products are disposed of improperly and seep into the ground, contaminating the water that fills the cracks and pore spaces in rocks and sediments beneath the earth's surface. In the context of Monterey Park, since 100\% of the water served to residents and businesses comes from groundwater, keeping this resource clean is of vital importance to the community.}}}
\end{refanswerenv}
\vspace{0.5em}
\noindent{\small\textcolor{usercolor}{\textbf{$\triangleright$\enspace User --- Turn~2}}}\hfill{\tiny\textcolor{gray}{Factoid, Summarization $\cdot$ Follow-up $\cdot$ ANSWERABLE}}
\begin{userbubble}
Can it be clean up?
\end{userbubble}
\noindent\hfill{\small\textcolor{agentcolor}{\textbf{Agent\enspace$\triangleleft$}}}
\begin{agentbubble}
Yes, contaminated groundwater can be cleaned up, but it can be a time-consuming and expensive process. Communities whose drinking water source has become contaminated have had to spend millions of dollars to remove contaminants from the water before it can be piped to homes and businesses. The cleanup is effective, but it can significantly affect the cost of providing the water to customers. It is far better to prevent contamination in the first place.
\end{agentbubble}
\noindent\hfill{\footnotesize\textcolor{doccolor}{\textbf{$\hookrightarrow$\enspace Retrieved documents~(2)}}}
\begin{docenv}
{\footnotesize\textcolor{doccolor}{\textbf{Doc~1}}}\hfill{\tiny\texttt{score:~23.123}}\ {\tiny\textcolor{gray}{\{'relevant': \{'9dcd142c-415e-4161-a82f-c8b14bfb36ac': \{'value': 'yes', 'timestamp': 1725439487\}\}\}}}\\[3pt]
{\small\textit{FAQs • Where do groundwater toxins come from?}}\\[2pt]
{\tiny\url{https://www.montereypark.ca.gov/FAQ.aspx?QID=118}}\\[4pt]
{\footnotesize FAQs • Where do groundwater toxins come from?
The water is imported for us by the San Gabriel Valley Municipal Water District, a public agency, and of which the City of Monterey Park is a member.	

4.
How much water does the Monterey Park community use?

Usage ranges from an average low of 7.3 million gallons per day (mgd) in February to an average high of 12 mgd in August. Average use per person is 100-110 gallons per day.	

5.
What is ground water?

Contrary to what a lot of people think, groundwater is not found in underground rivers or lakes. Rather, it is water that fills the cracks and pore spaces in rocks and sediments that lie beneath the surface of the earth, much the way water fills the open spaces and saturates a sponge or towel. 	

6.
Why is keeping our groundwater clean so important?

The San Gabriel Valley is unique in that it is one of the few natural underground reservoirs in the United States. Rain and snowmelt comes down from the San Gabriel Mountains through a series of dams and are then diverted to "spreading basins". Spreading basins are designated areas of land that are very permeable, such as sand and gravel areas. Water can then percolate into the ground. 100\% of the water served to residents and businesses in Monterey Park comes from groundwater. Thus, keeping our groundwater resource clean is vitally important to all of us. 	

7.
Where do groundwater toxins come from?

Many products we use around our homes everyday are considered toxic, such as household cleaners, flea powders, garden pesticides, paints, varnishes, motor oil, and antifreeze. These products are not only hazardous to our health, but to our environment. By disposing these harmful products improperly, we are contributing to the groundwater contamination in the San Gabriel Valley.	

8.
Can contaminated groundwater be cleaned up?

Yes, but it can be time-consuming and expensive. Communities whose drinking water source has become contaminated have had to spend millions of dollars to remove contaminants from the water before it can be piped to homes and businesses. The cleanup is effective, but it can significantly affect the cost of providing the water to customers.}
\end{docenv}
\begin{docenv}
{\footnotesize\textcolor{doccolor}{\textbf{Doc~2}}}\hfill{\tiny\texttt{score:~21.881}}\ {\tiny\textcolor{gray}{\{'relevant': \{'9dcd142c-415e-4161-a82f-c8b14bfb36ac': \{'value': 'yes', 'timestamp': 1725439552\}\}\}}}\\[3pt]
{\small\textit{FAQs • Where do groundwater toxins come from?}}\\[2pt]
{\tiny\url{https://www.montereypark.ca.gov/FAQ.aspx?QID=118}}\\[4pt]
{\footnotesize FAQs • Where do groundwater toxins come from?
By disposing these harmful products improperly, we are contributing to the groundwater contamination in the San Gabriel Valley.	

8.
Can contaminated groundwater be cleaned up?

Yes, but it can be time-consuming and expensive. Communities whose drinking water source has become contaminated have had to spend millions of dollars to remove contaminants from the water before it can be piped to homes and businesses. The cleanup is effective, but it can significantly affect the cost of providing the water to customers. It is far better to prevent contamination in the first place. What are the dos and do nots for groundwater?Do: Purchase only what you need Use products completely before disposing Recycle whatever you can (oil, batteries, etc.) Recycle surplus paintsDo not: Dump products onto the ground (they can contaminate soil and groundwater) Pour products or wastes into the street (they can contaminate oceans, lakes, and groundwater) Put hazardous or toxic products in the trash Pour products down the drainView proper disposal of Household Hazardous Wastes online. 

9.
Is our water fluoridated?

At this time fluoride is not added to the water. There are low levels of fluoride naturally occurring in the water. The range at which they occur is between 0.2 to 0.8 milligrams per liter (mg/L). The optimum fluoride level recommended to reduce dental caries (cavities) is a consistent .8 mg/L does. A bill signed by Governor Wilson in 1997 mandates drinking water fluoridation for larger water utilities, including Monterey Park, and requires the state to provide funding for the installation and operation of fluoridation equipment. The state is currently seeking funding for this program.	

10.
Does the City of Monterey Park own the Garvey Reservoir?

The Garvey Reservoir is not owned by the City of Monterey Park and the water it contains is not stored for city use.	

Live Edit

 

 

Coronavirus Updates

Agenda Center

Documents Prior to 2012

Teleconference Procedures

City Council Meetings Written Communications / Addendums

Written Communications / Addendums 2024

Written Communications / Addendums 2023}
\end{docenv}
\vspace{0.5em}
\noindent{\small\textcolor{usercolor}{\textbf{$\triangleright$\enspace User --- Turn~3}}}\hfill{\tiny\textcolor{gray}{Explanation $\cdot$ Clarification $\cdot$ ANSWERABLE}}
\begin{userbubble}
What is it?
\end{userbubble}
\noindent\hfill{\small\textcolor{agentcolor}{\textbf{Agent\enspace$\triangleleft$}}}
\begin{agentbubble}
Contrary to popular belief, groundwater is not found in underground rivers or lakes. Instead, water fills the cracks and pore spaces in rocks and sediments beneath the Earth's surface, similar to how water saturates a sponge or towel. Only 3.5 percent of the water on Earth is fresh water suitable for drinking, and most of that, 68 percent, is trapped in ice and glaciers. One-third of the freshwater is found in the ground, known as groundwater.
\end{agentbubble}
\noindent\hfill{\footnotesize\textcolor{doccolor}{\textbf{$\hookrightarrow$\enspace Retrieved documents~(2)}}}
\begin{docenv}
{\footnotesize\textcolor{doccolor}{\textbf{Doc~1}}}\hfill{\tiny\texttt{score:~24.299}}\ {\tiny\textcolor{gray}{\{'relevant': \{'9dcd142c-415e-4161-a82f-c8b14bfb36ac': \{'value': 'yes', 'timestamp': 1725440079\}\}\}}}\\[3pt]
{\small\textit{FAQs • Where do groundwater toxins come from?}}\\[2pt]
{\tiny\url{https://www.montereypark.ca.gov/FAQ.aspx?QID=118}}\\[4pt]
{\footnotesize FAQs • Where do groundwater toxins come from?
The water is imported for us by the San Gabriel Valley Municipal Water District, a public agency, and of which the City of Monterey Park is a member.	

4.
How much water does the Monterey Park community use?

Usage ranges from an average low of 7.3 million gallons per day (mgd) in February to an average high of 12 mgd in August. Average use per person is 100-110 gallons per day.	

5.
What is ground water?

Contrary to what a lot of people think, groundwater is not found in underground rivers or lakes. Rather, it is water that fills the cracks and pore spaces in rocks and sediments that lie beneath the surface of the earth, much the way water fills the open spaces and saturates a sponge or towel. 	

6.
Why is keeping our groundwater clean so important?

The San Gabriel Valley is unique in that it is one of the few natural underground reservoirs in the United States. Rain and snowmelt comes down from the San Gabriel Mountains through a series of dams and are then diverted to "spreading basins". Spreading basins are designated areas of land that are very permeable, such as sand and gravel areas. Water can then percolate into the ground. 100\% of the water served to residents and businesses in Monterey Park comes from groundwater. Thus, keeping our groundwater resource clean is vitally important to all of us. 	

7.
Where do groundwater toxins come from?

Many products we use around our homes everyday are considered toxic, such as household cleaners, flea powders, garden pesticides, paints, varnishes, motor oil, and antifreeze. These products are not only hazardous to our health, but to our environment. By disposing these harmful products improperly, we are contributing to the groundwater contamination in the San Gabriel Valley.	

8.
Can contaminated groundwater be cleaned up?

Yes, but it can be time-consuming and expensive. Communities whose drinking water source has become contaminated have had to spend millions of dollars to remove contaminants from the water before it can be piped to homes and businesses. The cleanup is effective, but it can significantly affect the cost of providing the water to customers.}
\end{docenv}
\begin{docenv}
{\footnotesize\textcolor{doccolor}{\textbf{Doc~2}}}\hfill{\tiny\texttt{score:~18.914}}\ {\tiny\textcolor{gray}{\{'relevant': \{'9dcd142c-415e-4161-a82f-c8b14bfb36ac': \{'value': 'yes', 'timestamp': 1725440074\}\}\}}}\\[3pt]
{\small\textit{How Much Water Is on Earth? | NASA Space Place – NASA Science for Kids}}\\[2pt]
{\tiny\url{https://spaceplace.nasa.gov/water/en/}}\\[4pt]
{\footnotesize How Much Water Is on Earth? | NASA Space Place – NASA Science for Kids


How Much Water Is on Earth? | NASA Space Place – NASA Science for Kids

water

Earth
Sun
Solar System
Universe
Science and Tech
Educators

How Much Water Is on Earth?

Click here to download this video (1920x1080, 37 MB, video/mp4).

How much water is on Earth?
Earth has a lot of water. It's one of the things that makes our big, blue planet so special. Learn more about Earth's water in this video! 

 Download the 8.5x11 poster!
 Download the 11x17 poster!
This poster summarizes the video and makes a great wall decoration!

 Click here to read a transcript of this video

How much water is on Earth?
Earth is special because it has so much water. It’s in the ground, forms big oceans, and it’s in our atmosphere.
It’s also in living things. Your body is made up mostly of water.
Water covers 71 percent of Earth’s surface. And almost all of it—96.5 percent—is salt water. 
Salt water has different kinds of salt in it, but it mostly has sodium chloride, the same salt we add to our food.
Just 3.5 percent of the water on Earth is fresh water we can drink. And most of that fresh water, 68 percent, is trapped in ice and glaciers.
A third of the fresh water is in the ground. We call it groundwater. The last two percent of fresh water is in the rivers, lakes, and streams.
A very small amount is in our atmosphere, where it exists as water vapor. That’s what clouds are made of.
If you took all the water on Earth and put it together it would be a ball 860 miles wide. That’s a lot of water!

Related Resources for Educators

Our World: ICESat-2 - What Is Ice?
Our World: Why Are Oceans Salty?
Real World: Global Cloud Observation Day

 	article last updated August 29, 2022
 

More about Earth

Make an Earth fan!

Earth's atmosphere

Why is thesky blue?

If you liked this, you may like:

How Does GPS Work?

What Is a Solar Eclipse?

Searching for Other Planets Like Ours

	 Games	 

	 Crafts	 

	 Activities	 

	 Videos	 

	 Glossary	 

	 Mystery	 

About Us
Privacy Policy
Image Use
Accessibility

Contact NASA Space Place

 Last Updated: May 6th, 2024 

 

More
Less

By Subject

Space
Sun
Earth
Solar System
People \& Technology
Parents \& Educators

By Type

Explore
Do
Play

More
Less

More about Earth

Make an Earth fan!

Earth's atmosphere

Why is thesky blue?

}
\end{docenv}
\begin{refanswerenv}
{\footnotesize\textit{\textcolor{gray}{\textbf{Reference answer:}\enspace Groundwater contamination is the presence of harmful substances in groundwater, which occurs when products like household cleaners, flea powders, garden pesticides, paints, varnishes, motor oil, and antifreeze are disposed of improperly and seep into the ground.
Cleanup is possible, but it can be time-consuming and expensive, often requiring communities to spend millions of dollars to remove contaminants before the water can be safely used. Preventing contamination in the first place is a more effective and affordable approach.}}}
\end{refanswerenv}
\vspace{0.5em}
\noindent{\small\textcolor{usercolor}{\textbf{$\triangleright$\enspace User --- Turn~4}}}\hfill{\tiny\textcolor{gray}{Summarization $\cdot$ Follow-up $\cdot$ ANSWERABLE}}
\begin{userbubble}
What are the suppliers for the city?
\end{userbubble}
\noindent\hfill{\small\textcolor{agentcolor}{\textbf{Agent\enspace$\triangleleft$}}}
\begin{agentbubble}
The City of Monterey Park Water System gets its water from local groundwater, which is pumped from 12 City-owned wells. These wells have a total capacity of 20 million gallons per day (mgd) and reach a depth of up to 1,800 feet. They are located near the Rio Hondo, outside the city limits, in the Main San Gabriel groundwater basin. Monterey Park's water is supplied by the San Gabriel Valley Municipal Water District, of which the city is a member. It's important to note that the City of Monterey Park does not own the Garvey Reservoir, and the water it contains is not stored for city use.
\end{agentbubble}
\noindent\hfill{\footnotesize\textcolor{doccolor}{\textbf{$\hookrightarrow$\enspace Retrieved documents~(2)}}}
\begin{docenv}
{\footnotesize\textcolor{doccolor}{\textbf{Doc~1}}}\hfill{\tiny\texttt{score:~28.664}}\ {\tiny\textcolor{gray}{\{'relevant': \{'9dcd142c-415e-4161-a82f-c8b14bfb36ac': \{'value': 'yes', 'timestamp': 1725442970\}\}\}}}\\[3pt]
{\small\textit{FAQs • Where do groundwater toxins come from?}}\\[2pt]
{\tiny\url{https://www.montereypark.ca.gov/FAQ.aspx?QID=118}}\\[4pt]
{\footnotesize FAQs • Where do groundwater toxins come from?
The water is imported for us by the San Gabriel Valley Municipal Water District, a public agency, and of which the City of Monterey Park is a member.	

4.
How much water does the Monterey Park community use?

Usage ranges from an average low of 7.3 million gallons per day (mgd) in February to an average high of 12 mgd in August. Average use per person is 100-110 gallons per day.	

5.
What is ground water?

Contrary to what a lot of people think, groundwater is not found in underground rivers or lakes. Rather, it is water that fills the cracks and pore spaces in rocks and sediments that lie beneath the surface of the earth, much the way water fills the open spaces and saturates a sponge or towel. 	

6.
Why is keeping our groundwater clean so important?

The San Gabriel Valley is unique in that it is one of the few natural underground reservoirs in the United States. Rain and snowmelt comes down from the San Gabriel Mountains through a series of dams and are then diverted to "spreading basins". Spreading basins are designated areas of land that are very permeable, such as sand and gravel areas. Water can then percolate into the ground. 100\% of the water served to residents and businesses in Monterey Park comes from groundwater. Thus, keeping our groundwater resource clean is vitally important to all of us. 	

7.
Where do groundwater toxins come from?

Many products we use around our homes everyday are considered toxic, such as household cleaners, flea powders, garden pesticides, paints, varnishes, motor oil, and antifreeze. These products are not only hazardous to our health, but to our environment. By disposing these harmful products improperly, we are contributing to the groundwater contamination in the San Gabriel Valley.	

8.
Can contaminated groundwater be cleaned up?

Yes, but it can be time-consuming and expensive. Communities whose drinking water source has become contaminated have had to spend millions of dollars to remove contaminants from the water before it can be piped to homes and businesses. The cleanup is effective, but it can significantly affect the cost of providing the water to customers.}
\end{docenv}
\begin{docenv}
{\footnotesize\textcolor{doccolor}{\textbf{Doc~2}}}\hfill{\tiny\texttt{score:~16.522}}\ {\tiny\textcolor{gray}{\{'relevant': \{'9dcd142c-415e-4161-a82f-c8b14bfb36ac': \{'value': 'yes', 'timestamp': 1725444071\}\}\}}}\\[3pt]
{\small\textit{FAQs • Monterey Park, CA • CivicEngage}}\\[2pt]
{\tiny\url{https://www.montereypark.ca.gov/faq.aspx}}\\[4pt]
{\footnotesize FAQs • Monterey Park, CA • CivicEngage
The City of Monterey Park Water System receives its water supply from one source - local groundwater. The water is produced by 12 City-owned wells with a total capacity of 20 million gallons per day (mgd). The wells reach a depth of up to 1,800 feet and are located in the vicinity of the Rio Hondo, which is outside the city limits and in the Main San Gabriel groundwater basin. Water from our deep wells is of excellent quality and requires little or no treatment. The Monterey Park Water System supplies an average of 10 mgd to its customers.On average, about 65\% of the water used each year is supplied from local rainfall; the other 35\% is imported from northern California and then percolated into the groundwater aquifers. The water is imported for us by the San Gabriel Valley Municipal Water District, a public agency, and of which the City of Monterey Park is a member.	

4.
How much water does the Monterey Park community use?

Usage ranges from an average low of 7.3 million gallons per day (mgd) in February to an average high of 12 mgd in August. Average use per person is 100-110 gallons per day.	

5.
What is ground water?

Contrary to what a lot of people think, groundwater is not found in underground rivers or lakes. Rather, it is water that fills the cracks and pore spaces in rocks and sediments that lie beneath the surface of the earth, much the way water fills the open spaces and saturates a sponge or towel. 	

6.
Why is keeping our groundwater clean so important?

The San Gabriel Valley is unique in that it is one of the few natural underground reservoirs in the United States. Rain and snowmelt comes down from the San Gabriel Mountains through a series of dams and are then diverted to "spreading basins". Spreading basins are designated areas of land that are very permeable, such as sand and gravel areas. Water can then percolate into the ground. 100\% of the water served to residents and businesses in Monterey Park comes from groundwater. Thus, keeping our groundwater resource clean is vitally important to all of us. 	

7.
Where do groundwater toxins come from?}
\end{docenv}
\begin{refanswerenv}
{\footnotesize\textit{\textcolor{gray}{\textbf{Reference answer:}\enspace The water for Monterey Park is imported by the San Gabriel Valley Municipal Water District, a public agency of which the City of Monterey Park is a member. The City of Monterey Park does not own the Garvey Reservoir, and the water it contains is not stored for city use.}}}
\end{refanswerenv}
\vspace{0.5em}
\noindent{\small\textcolor{usercolor}{\textbf{$\triangleright$\enspace User --- Turn~5}}}\hfill{\tiny\textcolor{gray}{Keyword $\cdot$ Clarification $\cdot$ ANSWERABLE}}
\begin{userbubble}
Pollution
\end{userbubble}
\noindent\hfill{\small\textcolor{agentcolor}{\textbf{Agent\enspace$\triangleleft$}}}
\begin{agentbubble}
Equipment, power plants that burn oil, gas, or coal, and industries that manufacture chemicals and other goods, along with other routine activities such as mowing the lawn, burning brush, heating with wood, and running errands in automobiles, all contribute to air pollution. Nearly everyone in any city around the world is exposed to air that has harmful levels of air pollution," said Anenberg, a member of NASA’s Health and Air Quality Applied Sciences team. For areas with little to no ground-based pollution monitoring, the data from Anenberg and colleagues may be the first that local officials have regarding the impacts of air pollution in their cities.
\end{agentbubble}
\noindent\hfill{\footnotesize\textcolor{doccolor}{\textbf{$\hookrightarrow$\enspace Retrieved documents~(2)}}}
\begin{docenv}
{\footnotesize\textcolor{doccolor}{\textbf{Doc~1}}}\hfill{\tiny\texttt{score:~11.068}}\ {\tiny\textcolor{gray}{\{'relevant': \{'9dcd142c-415e-4161-a82f-c8b14bfb36ac': \{'value': 'yes', 'timestamp': 1725445769\}\}\}}}\\[3pt]
{\small\textit{No Breathing Easy for City Dwellers: Particulates}}\\[2pt]
{\tiny\url{https://earthobservatory.nasa.gov/images/149580/no-breathing-easy-for-city-dwellers-particulates}}\\[4pt]
{\footnotesize No Breathing Easy for City Dwellers: Particulates
“Nearly everyone in any city around the world is exposed to air that has harmful levels of air pollution in it,” said Anenberg, a member of NASA’s Health and Air Quality Applied Sciences team.
For areas with little to no ground-based pollution monitoring, the data from Anenberg and colleagues may be the first that local officials have regarding the impacts of air pollution in their cities. The team plans to continue working with local stakeholders and governments to help them interpret the data.
“Even cities with ground monitors can benefit from this information because the monitors can be sparsely located and may not cover the entire urban area,” Anenberg said. The team next plans to apply their methodology and satellite data to examine trends over time for other pollutants, including ozone.
NASA Earth Observatory images by Joshua Stevens, using data courtesy of Southerland, V. A., et al. (2022). Story by Sara E. Pratt.

A new satellite-derived dataset links concentrations of fine particulate matter in air pollution with health outcomes in cities around the world.
Image of the Day for March 15, 2022

Instruments:
Aqua — MODIS
In situ Measurement
Model
OrbView-2 — SeaWiFS
Terra — MISR
Terra — MODIS

Appears in these Collections:
Air Quality
Applied Sciences

Image of the Day
Atmosphere
Life
Human Presence
Remote Sensing

View more Images of the Day:

Mar 14, 2022

Mar 16, 2022

References \& Resources

Anenberg, S.C., et al. (2022) Long-term trends in urban NO2 concentrations and associated paediatric asthma incidence: estimates from global datasets The Lancet Planetary Health, 6: e49–e58.
Hammer, M.S. et al. (2020) Global Estimates and Long-Term Trends of Fine Particulate Matter Concentrations (1998–2018) Environmental Science and Technology 54, 13, 7879–7890.
NASA Applied Sciences (2022, January 5) Child Asthma and Other Health Effects of Long-Term Urban Air Pollution. 
NASA Earth Observatory (2021, November 8) An Extra Air Pollution Burden.}
\end{docenv}
\begin{docenv}
{\footnotesize\textcolor{doccolor}{\textbf{Doc~2}}}\hfill{\tiny\texttt{score:~11.136}}\ {\tiny\textcolor{gray}{\{'relevant': \{'9dcd142c-415e-4161-a82f-c8b14bfb36ac': \{'value': 'yes', 'timestamp': 1725445766\}\}\}}}\\[3pt]
{\small\textit{Air - NYSDEC}}\\[2pt]
{\tiny\url{https://www.dec.ny.gov/environmental-protection/air-quality}}\\[4pt]
{\footnotesize Air - NYSDEC
Most pollutants come from:on- and off-road vehicles and power equipment,power plants that burn oil, gas or coal, andindustries that manufacture chemicals and other goods.Other routine activities, such as mowing the lawn, burning brush, heating with wood, and running errands in automobiles also create air pollution. Everyone can help significantly reduce pollution by following tips on Living the Green Life.

 Sign up with DEC Delivers to receive Air Mail! 
 

 The monthly air quality newsletter that includes public participation news and other air program issues in New York State.
 

Gov Delivery Form

You must have JavaScript enabled to use this form.

Email address

 
 This field is required.
 

Topic Id

 
- None -AdirondacksAir Quality (Air Mail! e-newsletter)Air Quality AlertsAmphibian Migrations and Road CrossingsBirdingBond Act UpdatesCatskillsChautauqua Lake News and UpdatesClimate ChangeContract Opportunities and M/WBE SupportDrinking Water Source Protection Program (DWSP2)Employment and InternshipsEncuentre Su Aventura (Find Your Adventure - Spanish Edition)Environmental JusticeEnvironmental Notice BulletinFive Rivers EECForests, Plants, and Land ConservationFreshwater Fishing and Boating (The Fishing Line)Global - Blank IdGrant OpportunitiesGreat Lakes Funding, News \& EventsGreen LivingHudson River AlmanacHudson River EstuaryHudson River Estuary Conservation and Land UseInvasive SpeciesOutdoor Discovery e-NewsletterPesticidesReinstein Woods EECSaltwater Fishing and BoatingShellfishingWastewater OperatorsWildlife, Fish, and Marine LifeCamping, Wildlife Viewing, and Hiking (Find Your Adventure)Hunting and TrappingSolid Waste and RecyclingMaking Waves

Leave this field blank

 

New York State's Air Pollution Control ProgramFor over 50 years, federal and state air pollution programs have included permits and technical requirements to control emission of criteria pollutants and toxic air contaminants from facilities, along with extensive measurement and monitoring of ambient pollutant levels.}
\end{docenv}
\begin{refanswerenv}
{\footnotesize\textit{\textcolor{gray}{\textbf{Reference answer:}\enspace Groundwater pollution or contamination refers to the presence of harmful substances, such as household cleaners, flea powders, garden pesticides, paints, varnishes, motor oil, and antifreeze, in groundwater. This occurs when these products are disposed of improperly and seep into the ground, contaminating the water that fills the cracks and pore spaces in rocks and sediments beneath the Earth's surface. Cleaning up contaminated groundwater can be a time-consuming and expensive process, so preventing contamination in the first place is crucial. The City of Monterey Park Water System obtains its water from local groundwater, specifically 12 City-owned wells with a total capacity of 20 million gallons per day (mgd), reaching depths of up to 1,800 feet and located near the Rio Hondo, outside the city limits, in the Main San Gabriel groundwater basin.}}}
\end{refanswerenv}
\vspace{0.5em}
\noindent{\small\textcolor{usercolor}{\textbf{$\triangleright$\enspace User --- Turn~6}}}\hfill{\tiny\textcolor{gray}{Non-Question $\cdot$ Clarification $\cdot$ ANSWERABLE}}
\begin{userbubble}
In water
\end{userbubble}
\noindent\hfill{\small\textcolor{agentcolor}{\textbf{Agent\enspace$\triangleleft$}}}
\begin{agentbubble}
Pollution occurs when water collects contaminants from pesticide use, mining, logging, and other sources and deposits them in water bodies. The Department of Pesticide Regulation (DPR) would be involved if pesticides were the pollutants of concern. The California Department of Forestry and Fire Protection (CDFFP) would comment on a TMDL if timber harvesting contributed to a water body's impairment. The Department of Fish and Wildlife would be involved if fish and wildlife were beneficial uses of water impacted in an impaired water body.
\end{agentbubble}
\noindent\hfill{\footnotesize\textcolor{doccolor}{\textbf{$\hookrightarrow$\enspace Retrieved documents~(2)}}}
\begin{docenv}
{\footnotesize\textcolor{doccolor}{\textbf{Doc~1}}}\hfill{\tiny\texttt{score:~14.910}}\ {\tiny\textcolor{gray}{\{'relevant': \{'9dcd142c-415e-4161-a82f-c8b14bfb36ac': \{'value': 'yes', 'timestamp': 1725446583\}\}\}}}\\[3pt]
{\small\textit{2001 Budget Analysis:State Water Resources }}\\[2pt]
{\tiny\url{https://lao.ca.gov/analysis\_2001/resources/res\_15\_3940.htm}}\\[4pt]
{\footnotesize 2001 Budget Analysis:State Water Resources 
nine regional water quality control boards, operating under SWRCB's oversight. The regions' lists and TMDLs are both submitted to SWRCB for its review and
approval.
A number of other state agencies assist in the development of TMDLs, depending on the sources of pollution and beneficial uses of a particular water body. For
example, the Department of Pesticide Regulation (DPR) would be involved if pesticides were the pollutants of concern, and the California Department of
Forestry and Fire Protection (CDFFP) would comment on a TMDL if timber harvesting was contributing to the impairment of a water body. The Department of
Fish and Game would be involved if fish and wildlife were beneficial uses of water that were impacted in an impaired water body.
The list of impaired water bodies and TMDLs are submitted to the U.S. Environmental Protection Agency (U.S. EPA) for its review and approval. If U.S. EPA
rejects a particular TMDL, it must establish the TMDL itself. In some cases, the development of TMDLs is carried out under a federal court settlement. These
TMDLs may be developed by U.S. EPA or delegated by U.S. EPA to the state to develop.
The TMDL Requirements Largely Ignored from 1972 to the Early 1990s. The requirement for TMDLs was largely ignored by federal and state water quality
agencies until the early 1990s. Until then, water quality regulation had been focused on controlling point sources of pollution through state and federal
permitting programs. Although water quality had been improving, over 40 percent of assessed water bodies nationwide remained too polluted for fishing or
swimming. It became apparent that the remaining water pollution problems—largely associated with nonpoint source pollution—required different solutions. The
activation of the TMDL program in the 1990s, in part instigated by lawsuits against the federal government to enforce the TMDL requirements, provides a
planning framework to address these problems.
Current State List of Impaired Water Bodies. The most recent list of impaired water bodies in the state (1998) lists 509 impaired water bodies for which 1,471
TMDLs have to be developed.}
\end{docenv}
\begin{docenv}
{\footnotesize\textcolor{doccolor}{\textbf{Doc~2}}}\hfill{\tiny\texttt{score:~16.711}}\ {\tiny\textcolor{gray}{\{'relevant': \{'9dcd142c-415e-4161-a82f-c8b14bfb36ac': \{'value': 'yes', 'timestamp': 1725446586\}\}\}}}\\[3pt]
{\small\textit{2001 Budget Analysis:State Water Resources }}\\[2pt]
{\tiny\url{https://lao.ca.gov/analysis\_2001/resources/res\_15\_3940.htm}}\\[4pt]
{\footnotesize 2001 Budget Analysis:State Water Resources 
In contrast, nonpoint source pollution is created when water picks up
contaminants from pesticide use, mining, logging, and other sources and deposits them in water bodies.
States are required to develop plans—called Total Maximum Daily Loads (TMDLs)—to meet water quality standards in the impaired water bodies that have been
listed. A TMDL is developed for each pollutant contributing to the impairment of a listed water body. The development of a TMDL involves the following
steps:

Determination of how much of a particular pollutant a water body can receive from all sources and still meet water quality standards. 
Determination of how much of the existing level of pollution must be reduced to meet water quality standards. 
Allocation of responsibility for reducing pollution to the acceptable level among the various contributing sources, including both point and nonpoint sources.
Development of a plan to implement the TMDL. While implementation plans are not currently a federal requirement, federal regulations effective later this
year will require these plans. 
Adoption of a regulation (a "basin plan amendment") to establish actions to implement a TMDL, such as revising permits to control waste discharges. 

When a listed water body later meets water quality standards, it is "de-listed." Neither state nor federal law specifies a time schedule for the development or
implementation of TMDLs or for when listed water bodies should become de-listed. Only three of the 18 impaired water bodies on the state's initial list from
1976 have been de-listed.
State Versus Federal Role. The state is primarily responsible for compiling the list of impaired water bodies and developing TMDLs. This is carried out by the
nine regional water quality control boards, operating under SWRCB's oversight. The regions' lists and TMDLs are both submitted to SWRCB for its review and
approval.
A number of other state agencies assist in the development of TMDLs, depending on the sources of pollution and beneficial uses of a particular water body. For
example, the Department of Pesticide Regulation (DPR) would be involved if pesticides were the pollutants of concern, and the California Department of}
\end{docenv}
\begin{refanswerenv}
{\footnotesize\textit{\textcolor{gray}{\textbf{Reference answer:}\enspace Water pollution refers to the contamination of water bodies such as rivers, lakes, and oceans, making them unsuitable for use. This can occur due to various reasons, including industrial waste disposal, sewage and wastewater discharge, oil spills, and agricultural runoff. In the case of Monterey Park, the primary source of water is local groundwater, which can also become contaminated due to improper disposal of harmful substances. As stated in the document, many household products such as cleaning supplies, pesticides, paints, and motor oil can contribute to groundwater contamination if disposed of improperly. Cleaning up contaminated groundwater can be a costly and time-consuming process, emphasizing the importance of preventing contamination in the first place.}}}
\end{refanswerenv}
\vspace{0.5em}
\noindent{\small\textcolor{usercolor}{\textbf{$\triangleright$\enspace User --- Turn~7}}}\hfill{\tiny\textcolor{gray}{Summarization $\cdot$ Follow-up $\cdot$ ANSWERABLE}}
\begin{userbubble}
Sollutions
\end{userbubble}
\noindent\hfill{\small\textcolor{agentcolor}{\textbf{Agent\enspace$\triangleleft$}}}
\begin{agentbubble}
To reduce air pollution, here are some suggestions:

1. Turn off the lights when you leave a room.
2. Replace energy-hungry incandescent lights with compact fluorescent light bulbs.
3. Ask your energy supplier for a home audit and inquire about alternative energy solutions like solar or wind.
4. Opt for a fan instead of air conditioning.
5. Use a programmable thermostat and set it to 78°F in the summer and 68°F in the winter.
6. Install low-flow shower heads.
7. Recycle paper, plastic, metals, and organic materials.
8. Use an EPA-approved wood-burning stove or replace insert.
9. Don’t use your wood stove or replace on days with unhealthy air.
10. Don’t heat your home with a gas stove.
11. Use a surge protector for multiple appliances and turn it off when products are not in use.
12. Add insulation to your home.
13. Wash laundry in cold water and line dry.
To reduce water pollution, you can:
1. Pick up after your pet. Pet waste can contain harmful bacteria and should be disposed of properly.
2. Save water by sweeping your driveway instead of using the hose, and consider washing your car at a facility with advanced filtration systems to recycle water.
3. Consider going organic.
\end{agentbubble}
\noindent\hfill{\footnotesize\textcolor{doccolor}{\textbf{$\hookrightarrow$\enspace Retrieved documents~(2)}}}
\begin{docenv}
{\footnotesize\textcolor{doccolor}{\textbf{Doc~1}}}\hfill{\tiny\texttt{score:~15.930}}\ {\tiny\textcolor{gray}{\{'relevant': \{'9dcd142c-415e-4161-a82f-c8b14bfb36ac': \{'value': 'yes', 'timestamp': 1725447123\}\}\}}}\\[3pt]
{\small\textit{Simple Solutions to Help Reduce Air Pollution | California Air Resources Board}}\\[2pt]
{\tiny\url{https://ww2.arb.ca.gov/resources/fact-sheets/simple-solutions-help-reduce-air-pollution}}\\[4pt]
{\footnotesize Simple Solutions to Help Reduce Air Pollution | California Air Resources Board
The suggestions below will help reduce exposure in your home:Turn the lights off when you leave a room.Replace energy-hungry incandescent lights with compact fluorescent light bulbs.Ask your energy supplier for a home audit and inquire about alternative energy solutions like solar or wind.Opt for a fan instead of air conditioning.Use a programmable thermostat and set it to 78°F in the summer and 68°F in the winter.Install low- ow shower heads.Recycle paper, plastic, metals and organic materials.Use an EPA-approved wood burning stove or replace insert.Don’t use your wood stove or replace on days with unhealthy air.Don’t heat your home with a gas stove.Use a surge protector for multiple appliances and turn it off when products are not in use.Add insulation to your home.Wash laundry in cold water and line dry.When ready to replace, look for energy star appliances.Use a propane or natural gas barbecue rather than a charcoal one.Microwave or use a toaster oven for small meals.Have your gas appliances and heater regularly inspected and maintained.Use washable dishes, utensils and fabric napkins rather than disposable dinnerware.Choose products that use recycled materials.Eat locally, shop at farmers markets and buy organic products.Buy products from sustainable sources such as bamboo and hemp.Use durable reusable grocery bags and keep them in your car so you’re never caught off guard.Paint with a brush instead of a sprayer.Store all solvents in airtight containers.Use an electric or push lawn mower.Use a rake or broom instead of a leaf blower.Use water-based cleaning products that are labeled ‘zero VOC’.Insulate your water heater and any accessible hot water pipes.Eliminate use of toxic chemicals at home; opt for natural substitutes.Plant a tree! They filter the air and provide shade.Let your elected representatives know you support action for cleaner air.At WorkThere are multiple ways of reducing consumption at the workplace.}
\end{docenv}
\begin{docenv}
{\footnotesize\textcolor{doccolor}{\textbf{Doc~2}}}\hfill{\tiny\texttt{score:~15.465}}\ {\tiny\textcolor{gray}{\{'relevant': \{'9dcd142c-415e-4161-a82f-c8b14bfb36ac': \{'value': 'yes', 'timestamp': 1725447126\}\}\}}}\\[3pt]
{\small\textit{Rain Brings Renewal, But Also Toxins: How Californians Can Keep Their Water Pollutant-Free This Summer | Caltrans}}\\[2pt]
{\tiny\url{https://dot.ca.gov/news-releases/news-release-2023-022}}\\[4pt]
{\footnotesize Rain Brings Renewal, But Also Toxins: How Californians Can Keep Their Water Pollutant-Free This Summer | Caltrans
California residents unknowingly contribute to water pollution through their daily activities. The use of toxic chemical pesticides or neglecting vehicular maintenance, can further pollute our water supply. These substances inevitably end up in our waterways, contaminating local lakes, rivers, streams and the ocean. Thankfully, there are simple steps all of us can take to reduce stormwater pollution and help keep our water clean.
Californians can positively impact the environment by taking necessary steps to minimize contributions to water pollution. The “Let’s Change This to That” campaign via CleanWaterCA.com provides a list of tips to help support clean water from summer to winter and all year round. They include:
Pick up after your pet. A common misconception is that pet waste is a natural fertilizer and can be left to decompose in the environment. Not so! Pet waste can contain harmful bacteria such as E. coli and fecal coliform, which can carry diseases that affect wildlife and contaminate waterways when not properly disposed of.
Save water. Reduce water pollution by sweeping your driveway instead of using the hose and, if you can, by washing your car at a facility. The United States Environmental Protection Agency warns that washing cars and driveways at home can cause pollution as metals and fluids are washed off and enter storm drains, collecting more pollutants along the way and ending up in our waterways. If you can, choose car wash facilities that use advanced filtration systems to recycle water used for washing cars.
Go organic. This summer consider using natural pest control methods instead of harmful pesticides. Pesticides made from highly toxic chemicals can have long-lasting effects on the health of our water and the environment. By choosing organic alternatives, like companion planting or using beneficial insects like ladybugs, you can make a difference.
Give your vehicle a tune-up. Vehicle fluids, oils and brake dust are made up of toxic heavy metals that end up on roadways and highways if we don’t provide proper maintenance.}
\end{docenv}
\begin{refanswerenv}
{\footnotesize\textit{\textcolor{gray}{\textbf{Reference answer:}\enspace To prevent groundwater contamination, it is recommended to follow these guidelines:
* Purchase only what you need
* Use products completely before disposing
* Recycle whatever you can (oil, batteries, etc.)
* Recycle surplus paints
On the other hand, the following practices should be avoided:
* Dumping products onto the ground (they can contaminate soil and groundwater)
* Pouring products or wastes into the street (they can contaminate oceans, lakes, and groundwater)
* Putting hazardous or toxic products in the trash
* Pouring products down the drain
For proper disposal of Household Hazardous Wastes, you can view the instructions online.
Regarding water pollution, it is caused by various activities such as pesticide use, mining, logging, and other sources. To address this issue, relevant departments like the Department of Pesticide Regulation (DPR), California Department of Forestry and Fire Protection (CDFFP), and the Department of Fish and Wildlife should be involved in implementing appropriate measures and regulations.}}}
\end{refanswerenv}
\end{document}
